% Style-independent scientific content. Official ICLR 2027 wrapper supplied separately.
% Controlled-study values must be supplied by the shared audited result source.
\providecommand{\ControlledShiftAbstractResult}{The first preregistered controlled-shift primary made 38 VERA deployments with 0 violations but failed its usefulness gate because the 95\% lower bound was 15.1\%. We then locked an independent, non-pooling follow-up before seeds 109--172. At the higher preregistered 8,000-observation budget, validation selection violated 24/186 deployed decisions while VERA deployed 59 edits with 0 shifted-contract violations, retained 59/183 oracle-safe opportunities (32.2\%, 95\% CI [27.5\%, 36.8\%]), and passed every registered follow-up gate.}
\providecommand{\ControlledPrimaryResult}{The first controlled-shift primary remains a disclosed mixed result: validation selection violated 19/186 deployed decisions, whereas the VERA vector envelope deployed 38 edits with 0 violations; safety, paired reduction, and vector/common advantage passed, but usefulness failed because safe retention was 38/187 (20.3\%) with 95\% seed-bootstrap CI [15.1\%, 25.5\%]. The independent post-failure follow-up is not pooled with that result, but it passed as a new locked test: validation selection violated 24/186 deployed decisions, VERA deployed 59 edits with 0 violations, and the paired sign test had 22 favorable seed clusters, 0 adverse, and $p=4.77e-07$.}
\providecommand{\ControlledSafetyResult}{In the independent follow-up, the rotating-sentinel safety gate passed with 0/64 false acceptances and a one-sided 95\% Clopper--Pearson upper bound of 4.6\%. Across all follow-up VERA vector deployments, shifted-law violations were 0/59.}
\providecommand{\ControlledRetentionResult}{Follow-up safe retention was 59/183 (32.2\%) with 95\% seed-bootstrap CI [27.5\%, 36.8\%], clearing the registered 20\% lower-bound gate. The common-radius rule retained 11/183 (6.0\%); the vector/common ratio was 5.36 with 95\% interval [3.5, 10.4], so the vector-advantage gate also passed.}
\providecommand{\ControlledAllocationResult}{The follow-up evidence-size sweep shows the effect of certification budget: at $\Gamma=1.1$, the targeted vector rule retained 2/183, 22/183, 51/183, and 59/183 oracle-safe opportunities at budgets 1,000, 2,000, 4,000, and 8,000, respectively, with zero deployed violations in every row. The 8,000-observation row was locked as the independent follow-up primary before seeds 109--172 were run.}
\providecommand{\ControlledEvidenceResult}{The follow-up geometry remains narrow but useful: among 59 vector deployments, the median common radius was 1.040, with range [1.002, 1.183]. The most frequent limiting coordinate was target::environment=0 (39 of 59 deployments), followed by source::0 (14 of 59).}
\providecommand{\ControlledGaitResult}{The follow-up did not deploy uniformly across domains: VERA vector abstained on Bios and CivilComments, deployed 14/64 GaitPDB seed decisions and 45/64 Waterbirds seed decisions, and had zero shifted-contract violations on every deployment. This is the intended abstain-or-edit behavior under the supported-envelope contract.}
\providecommand{\ControlledHeldoutResult}{The held-out boosted-tree stress remains outside the formal registered attacker guarantee: it found 5 violations among 59 follow-up vector deployments (8.5\% unsafe), so the paper reports it as portfolio-scope stress rather than as a coverage claim.}
\providecommand{\ControlledNegativeResults}{The first 64-seed primary's failed usefulness gate remains mandatory negative evidence and is not rescued, pooled, or replaced by the follow-up. The independent follow-up is reported as a separately locked post-failure confirmation at the preregistered higher evidence budget; secondary threshold-stress, held-out-attacker, and dataset-specific rows remain ineligible as primary substitutes.}
\providecommand{\ControlledMainResultTable}{\input{../aaai2027_template/AuthorKit27/vera_followup_results_table}}
\providecommand{\ControlledMainResultFigure}{\begin{figure}[t]\centering\includegraphics[width=\columnwidth]{../figures/vera_controlled_shift_followup_budget_curve.pdf}\caption{Independent follow-up evidence-size sensitivity at $\Gamma=1.1$. The circled point is the locked follow-up primary vector rule (budget 8,000, targeted allocation), which clears the registered usefulness gate with zero shifted-contract violations. Lower budgets and uniform allocation are reported as sensitivity analyses.}\label{fig:controlled-followup-budget}\end{figure}}
\providecommand{\VERAMethodOverviewPath}{../figures/vera_method_overview.pdf}

\begin{abstract}
Representation editing changes what downstream models can recover from a
frozen representation, but an edit that looks acceptable on IID validation
data may damage its target task or expose a source concept after deployment
reweighting. We introduce VERA, a decision layer for a registered edit and a
declared supported shift. VERA audits paired target harm and a heterogeneous
portfolio of freshly retrained attackers, returning a simultaneous vector
envelope of environment- and source-conditional density-ratio budgets, its
limiting common radius, or \textsc{Abstain}. Finite-family testing follows
Learn Then Test; our contribution is the representation-intervention contract
that couples paired harm, balanced recovery, a continuum of groupwise shifts,
and an explicit support boundary. We prove finite-sample validity throughout
the reported envelope, derive prospective evidence allocations for additive
multi-cell contracts, and show that no certification-data-only protocol can certify a
required cell absent from its support without extra assumptions. In a prior IID
study, validation-only selection violated 35 of 128 contracts, compared with 1
of 128 for fixed-profile certification; this does not establish robustness for
non-IID shifts. \ControlledShiftAbstractResult
\end{abstract}

\section{Representation Edits Need Deployment Decisions}
\label{iclr:sec:intro}

Methods such as INLP, R-LACE, LEACE, kernel erasure, TaCo, SPLINCE, and MANCE++
construct interventions intended to remove a concept while preserving useful
information
\cite{ravfogel2020null,ravfogel2022linear,belrose2023leace,ravfogel2022kernel,jourdan2023taco,holstege2025splince,avitan2026mance}.
Their outputs are candidates, not complete deployment decisions. A target label
may be correlated with the erased concept, an eraser's native probe may miss
information recoverable by another hypothesis class, and group prevalence may
change after validation. A useful audit must therefore answer more than whether
one probe's average accuracy fell.

We ask: \emph{should this registered representation edit be deployed under this
declared shift?} VERA fits no new eraser. It evaluates edits produced without
certification data, retrains downstream target and leakage models after every
edit, and either certifies one edit within a supported shift class or abstains.
Its target quantity is incremental paired harm relative to the identity
representation. Its leakage quantity is equal-weight source-class recall for
each registered attacker, preventing source prevalence from creating apparent
erasure. Its output is anisotropic: it reports which environment, source class,
or attacker limits a deployment statement instead of hiding all constraints in
one average score.

We make four contributions. First, we define a support-aware certificate for a
representation intervention that couples paired target harm with balanced
recovery by a retrained attacker portfolio. Second, we obtain simultaneous
finite-sample validity over registered edits, contracts, and a continuum of
supported reweighting profiles, and derive prospective evidence allocations
for additive multi-cell contracts. Third, a two-world impossibility theorem
shows why the certificate must collapse on required support absent from
certification. Fourth, we conduct a preregistered controlled cross-modal study that
separates eraser construction from nine matched deployment rules and measures
unsafe deployment, useful retention, envelope geometry, and abstention.

The scope is deliberately narrower than universal erasure. Passing VERA does
not prove that all measurable information about a concept is gone, that a
representation is fair or private, or that an unobserved deployment belongs to
the declared ambiguity class. It states that one registered intervention meets
specific target and attacker contracts over a supported reweighting set.

\begin{figure*}[t]
\centering
\includegraphics[width=\textwidth]{\VERAMethodOverviewPath}
\caption{\textbf{VERA allocates evidence, audits one profile, and exposes its
support boundary.} (A) Independent design-fold margins determine disjoint
certification-stream counts before certification outcomes are observed. Paired
identity/edit predictions measure target harm, while one source-class draw
evaluates every registered fresh attacker. The primary controlled study uses
the preregistered square-score allocator; the disjoint-stream additive
comparison is a separately, prospectively specified secondary analysis. (B)
The candidate IUT decides the registered profile, while simultaneous schematic
upper-risk curves provide an inspectable envelope and limiting common radius.
(C) The guarantee cannot cross new support: Camelyon17 certification centers
are 0, 1, 3, and 4, whereas external center 2 is absent from that support, so
VERA abstains without an additional cross-center assumption. This is an
identification boundary, not evidence of clinical harm.}
\label{iclr:fig:overview}
\end{figure*}

\section{Closest Work}
\label{iclr:sec:related}

\paragraph{Erasure and probing.}
Linear erasure includes iterative nullspace projection, adversarial projection,
and closed-form LEACE
\cite{ravfogel2020null,ravfogel2022linear,belrose2023leace}. Nonlinear and
structure-aware approaches broaden the candidate class
\cite{ravfogel2022kernel,chowdhury2023kram,shao2023unaligned,jourdan2023taco,holstege2025splince,avitan2026mance}.
These methods differ in what interventions they can construct; VERA instead
compares their outputs under one downstream audit. Probe conclusions depend on
the hypothesis class and retraining protocol
\cite{hewitt2019control,pimentel2020probing,voita2020mdl,elazar2021amnesic},
which motivates a registered heterogeneous portfolio rather than the eraser's
native linear objective alone. Fundamental utility limits and residual
nonlinear recoverability also caution against treating optimization success as
a deployment guarantee \cite{chowdhury2025limits,akbari2025obliviator}.
FARE and MMD-B-Fair provide certificates for restricted fair-representation
objectives \cite{jovanovic2023fare,deka2023mmdbfair}; VERA does not supersede
those objectives or certify universal guardedness.

\paragraph{Risk control and selection.}
Risk-controlling prediction sets, Learn Then Test (LTT), Conformal Risk Control,
and Pareto Testing provide the finite-family testing foundation
\cite{bates2021rcps,angelopoulos2025ltt,angelopoulos2024crc,laufer2023pareto}.
Prompt Risk Control is a close application template for selecting prompts under
multiple risks and estimated shift \cite{zollo2024prompt}; selective, joint,
and group-conditional certificates further show that finite-family, multi-risk,
selective, and grouped acceptance are not themselves new
\cite{xu2025scrc,yu2026joint,huang2026conditional}. VERA inherits candidate-wise testing.
Its proposed delta is the audited population object: paired intervention harm,
source-balanced recovery by retrained attackers, a vector of supported shift
budgets, and an explicit missing-support outcome.

\paragraph{Shift and abstention.}
Known weights and robust validation can control risk under stated covariate or
ambiguity models
\cite{tibshirani2019covariate,almeida2025hprc,zecchin2025wcrc,cauchois2024robust,ai2024finegrained}.
Robust evaluation, CVaR concentration, and distributional fairness
certificates are close technical and application neighbors
\cite{najafi2025federated,thomas2019cvar,kang2022certifying,ehyaei2025fragile}.
Group DRO and distribution-shift benchmarks motivate conditional rather than
globally averaged evaluation \cite{sagawa2020groupdro,koh2021wilds}, while
domain-adaptation impossibility and robustness-evaluation work motivate the
explicit support boundary \cite{bendavid2010impossibility,subbaswamy2021evaluation}.
Selective classification rejects individual predictions
\cite{elyaniv2010foundations,geifman2017selective,geifman2019selectivenet}; VERA
rejects an entire intervention before deployment. Three distinctions matter in
practice. IID certification is only the radius-one slice of the proposed
object; scalar pooling can accept while one target environment or attacker
fails; and no confidence interval can identify an omitted required cell.

\section{Representation-Intervention Contract}
\label{iclr:sec:contract}

Let $Z\in\mathbb R^d$ be a frozen representation, $Y$ a target, $S\in\{0,1\}$
a source concept, and $G$ an observed environment. A registered edit
$e\in\mathcal E$ maps $Z$ to $Z^e$ and is fitted without certification data.
Fresh target learners $f_0$ and $f_e$ are fitted for identity and edited
representations. On certification example $i$, define paired harm
\begin{equation}
H_{e,i}=\mathbf 1\{f_e(Z_i^e)\ne Y_i\}-
\mathbf 1\{f_0(Z_i)\ne Y_i\}\in\{-1,0,1\}.
\label{iclr:eq:harm}
\end{equation}
Pairing subtracts baseline difficulty on the same example. For attacker
$a\in\mathcal A$, retrained after edit $e$, let
$L_{e,a,i}=\mathbf 1\{a_e(Z_i^e)=S_i\}$.

For a reference law $P$ and cap $\Gamma\ge1$, define
\begin{equation}
\mathcal Q_\Gamma(P)=\left\{Q\ll P:
0\le\frac{\mathrm dQ}{\mathrm dP}\le\Gamma\right\},\qquad
\rho_{\Gamma,P}(V)=\sup_{Q\in\mathcal Q_\Gamma(P)}\mathbb E_Q[V].
\label{iclr:eq:robust-risk}
\end{equation}
For environment $g$, the target contract is
$\rho_{\gamma_g,P_g}(H_e)\le\tau$. Leakage uses the standardized robust
balanced accuracy
\begin{equation}
B_{e,a}(\eta_0,\eta_1)=\frac12\sum_{s\in\{0,1\}}
\rho_{\eta_s,P_s}(L_{e,a})\le\lambda.
\label{iclr:eq:balanced}
\end{equation}
Equal source-class weighting makes the contract invariant to deployment source
prevalence; conditioning harm on $G$ prevents a favorable environment mixture
from hiding damage. The target and leakage conditionals need not form one
factorized joint model. Each declared conditional law must satisfy its own cap.

The contract is fixed before certification: candidate family, target learner,
attacker portfolio, thresholds, support cells, and deployment profile are all
registered. Identity is the paired comparator, not silently added as another
selectable action. If identity can be selected, it belongs in the corrected
finite family.

\section{VERA Envelope and Decision}
\label{iclr:sec:method}

The bounded-reweighting risk has the CVaR representation
\begin{equation}
\rho_{\Gamma,P}(V)=\inf_{t}
\{t+\Gamma\mathbb E_P[(V-t)_+]\}.
\label{iclr:eq:cvar}
\end{equation}
A simultaneous distribution-function band therefore gives an upper risk curve
for every $\Gamma\ge1$. The audit variables here admit exact simplifications.
For Bernoulli leakage with mean $p$,
$\rho_{\Gamma,P}(L)=\min(1,\Gamma p)$. For paired harm with probabilities
$(p_-,p_0,p_+)$ on $(-1,0,1)$, write $a=\Gamma p_+$ and
$b=\Gamma(1-p_-)$. Then
\begin{equation}
\rho_{\Gamma,P}(H)=
\begin{cases}
1,&a\ge1,\\
a,&a<1\le b,\\
a+b-1,&b<1.
\end{cases}
\label{iclr:eq:ternary}
\end{equation}
VERA substitutes one-sided exact binomial bounds, preserving the probability
simplex for the two harm tails and splitting component error across source
classes for balanced leakage.

At one declared profile, each candidate is an intersection--union test (IUT):
every target and leakage component must clear its threshold using candidate
error $\delta/|\mathcal E|$. If an unsafe candidate is accepted, at least one
fixed violated component must undercover; intersecting additional components
does not require another union bound. A separate simultaneous construction
covers every curve for every candidate and contract because the curves remain
available for post-certification inspection.

Let $U^T_{e,g}(\gamma_g)$ be the simultaneous target upper curve and
$U^L_{e,a,s}(\eta_s)$ the source-class attacker curve. VERA reports
\begin{equation}
\widehat{\mathcal E}_e=\left\{(\boldsymbol\gamma,\boldsymbol\eta):
U^T_{e,g}(\gamma_g)\le\tau\ \forall g,\quad
\frac12\sum_sU^L_{e,a,s}(\eta_s)\le\lambda\ \forall a\right\}.
\label{iclr:eq:envelope}
\end{equation}
An unsupported required coordinate has radius zero. The common radius is the
largest diagonal profile inside this set, right-censored at the registered
computational cap. The full vector reveals which constraint limits the edit;
the diagonal is only a summary.

\begin{theorem}[Supported intervention certificate]
\label{iclr:thm:certificate}
Conditional on construction and all registered audit objects, the fixed-profile
IUT deploys no contract-violating candidate with probability at least
$1-\delta$. Under simultaneous candidate--contract coverage, the same event
ensures that every profile in every reported $\widehat{\mathcal E}_e$ satisfies all registered
target and balanced-leakage contracts.
\end{theorem}

The proof is inherited finite-family testing applied to intervention-specific
components. The continuum does not incur an uncountable union bound: one
simultaneous distribution-function event implies the entire monotone robust-risk
curve. Environment mixtures remain controlled by linearity, and source
prevalence is absent from the standardized leakage functional. Utility may be
used to choose among candidates already contained in the same coverage event.

\section{Guarantees and Evidence Design}
\label{iclr:sec:evidence}

For a bounded component with range $R$, profile cap $\Gamma_0$, and population
margin $m>0$, a DKW construction has power at least $1-\beta/J$ when
\begin{equation}
n\ge \frac{\Gamma_0^2R^2}{2m^2}
\left(\sqrt{\log\frac{2}{\alpha}}+
\sqrt{\log\frac{2J}{\beta}}\right)^2.
\label{iclr:eq:sample-size}
\end{equation}
This makes abstention an evidence diagnosis rather than an opaque failure.

For additive contracts, let $a_{jc}$ collect cell $c$'s range, budget,
confidence, and margin factors in contract $j$. Under the registered separate
target-environment and source-class streams, VERA solves
\begin{equation}
\min_{\boldsymbol n}\max_j\sum_c\frac{a_{jc}}{\sqrt{n_c}}
\quad\text{s.t.}\quad \sum_cn_c=N,\quad n_c\ge n_c^{\min}.
\label{iclr:eq:allocation}
\end{equation}
The program is convex. For one additive contract with inactive floors,
$n_c\propto a_{1c}^{2/3}$; when each contract uses one cell, the solution
reduces to square-score allocation after shared-cell maxima are taken. Margins
are estimated on an independent design fold for its max-min-margin edit, so the
allocation targets one registered power surrogate without affecting
certificate validity for the whole candidate family.

The matching statement is local, not universal. Within its stated independent
Bernoulli comparison experiment, the exact-KL lower bound needs no interior
restriction; the familiar $\Gamma^2/m^2$ scaling requires safe and unsafe
probabilities inside a fixed compact subset of $(0,1)$. Allocation
optimality is restricted to the DKW surrogate and the registered separate-stream
architecture. It does not cover a fully joint environment--source design that
reuses every observation across target and leakage contracts.

\begin{theorem}[Unsupported-cell impossibility]
\label{iclr:thm:impossibility}
If a required cell has zero certification probability and the model class
contains safe and unsafe worlds that induce identical protocol observations but
disagree on that cell's contract, any certification-data-only protocol with uniform
false-acceptance control $\delta$ accepts in the safe world with probability at
most $\delta$ whenever deployment may concentrate on that cell.
\end{theorem}

Identical observations force identical acceptance probabilities in the two
worlds. Safety in the unsafe world therefore forces abstention in the safe one.
More optimization cannot remove this boundary; the remedies are new support or
an explicit structural assumption linking observed and omitted cells.

\section{Controlled Cross-Modal Study}
\label{iclr:sec:protocol}

The prospective study crosses four supported datasets, five eraser families,
and 64 fresh seeds (45--108), yielding 1,280 dataset--method runs and 3,072
candidate-edit evaluations. Waterbirds provides vision; CivilComments and Bios
provide text; GaitPDB provides time series. Camelyon17 is reserved for the
unsupported hospital boundary because deployment center 2 is absent from
certification support
\cite{koh2021wilds,sagawa2020groupdro,borkan2019civilcomments,dearteaga2019bios}.
The supported study concerns exact reweightings of finite benchmark reference
laws, not new people, images, comments, or recordings.

The candidate frontier has 12 edit configurations per dataset--seed block:
INLP ranks $\{1,2,4,8\}$, R-LACE ranks $\{1,4\}$, one LEACE solution, TaCo
removal counts $\{1,2,3,5\}$, and one MANCE++ configuration. Every method uses a
pinned upstream entry point under one frozen-feature protocol. This is entry-
point parity, not reproduction of every upstream paper's full pipeline or
eraser state of the art; no proxy implementation contributes a candidate or
result. After each edit, the study retrains a logistic target
learner and four leakage attackers:
logistic regression, RBF--Nystroem logistic regression, random forest, and a
two-layer MLP. A boosted tree is fitted on the training partition but withheld
from edit construction, selection, and certification.

The official training split is divided into eraser-training and construction
partitions. Pinned edit adapters use them as required, while fresh target and
attacker models fit only the eraser-training partition. A disjoint design fold chooses the targeted
supported cell and evidence allocation. Independent certification draws form
confidence bounds. Primary exact shifted-law labels use expectations under a
known finite reweighting, with a separate 50,000-draw Monte Carlo replay. Every
receipt records support, probability, profile, split, upstream-code, and audit-
array hashes. The primary requests global cap $\Gamma=1.1$ and uses the exact
induced conditional profile, with $\delta=0.05$ and 4,000 contract observations.
The fixed $(\tau,\lambda)$ values are $(0.40,0.70)$ for
Bios, $(0.075,0.80)$ for CivilComments, $(0.20,0.55)$ for GaitPDB, and
$(0.10,0.90)$ for Waterbirds. They come from an earlier design-only seed block
and are benchmark operating points, not normative safety tolerances. The
targeted primary gives every registered evidence cell at least 15\% of the
budget; uniform allocation is matched. Secondary profiles vary the requested
global cap over $\Gamma\in\{1.1,1.25,1.5\}$ and
$N\in\{1000,2000,4000,8000\}$. Common radii are right-censored at
$\Gamma_{\max}=8$.
The preregistration sets the inversion cap to 8, whereas an earlier analysis
implementation used cap 4. Receipt generation and fixed-profile certification
are unaffected. We treat cap 8 as the prespecified analysis, retain the
unchanged cap-4 output as a sensitivity, require agreement with an independent
cap-8 replay, and report every resulting decision or gate difference.

The 64-seed sample size was fixed from the disjoint seeds 13--44 design block.
Its 14 favorable, zero adverse, and 18 tied seed-level differences gave
Jeffreys-smoothed exact-sign-test power above 0.999 at 64 seeds. Sixty-four also
provides 16 rotating sentinel decisions per dataset and exceeds the exact
59-seed minimum needed for a zero-event one-sided 95\% upper bound of 0.05.

Nine rules share candidates, audit variables, thresholds, shifted laws, and
seed clusters: always deploy, validation point selection, IID LTT, robust point
selection, a generic scalar robust certificate, VERA fixed-profile IUT, the
vector envelope, the common-radius restriction, and an exact shifted-law
opportunity oracle. The scalar rule averages normalized robust upper bounds, whereas VERA
requires each coordinate to pass. The independent unit is the entire seed
cluster, not datasets, candidates, attackers, or profiles within a seed.

A predefined supporting stress recomputes always deploy, validation point
selection, and the vector envelope at
$\kappa\in\{0.75,1,1.25\}$ using
$\tau_d(\kappa)=\kappa\tau_d$ and
$\lambda_d(\kappa)=0.5+\kappa(\lambda_d-0.5)$. It reports all 36 cells and
cannot change a primary gate. Its main-track supporting condition requires
both naive rules to reach 20\% violations somewhere at the registered
$\kappa=1$, VERA to remain at or below 5\% in every dataset--severity cell,
and the registered safety and usefulness gates to pass.

The primary efficacy endpoint is the within-seed reduction in violations from
point selection to the vector rule, tested by a two-sided exact sign test on
nonzero seed differences. The rotating-dataset safety sentinel uses one
prespecified decision per seed and a one-sided exact 95\% Clopper--Pearson upper
bound; any observed VERA false acceptance fails. Useful abstention requires a
95\% percentile-bootstrap lower bound of at least 20\% safe retention from
20,000 whole-seed resamples with fixed seed 2027071601, and vector retention
must be at least twice common-radius retention. These are an IUT: every gate
must pass. Per-dataset tests are secondary and Holm adjusted.

\section{Deployment Safety and Useful Abstention}
\label{iclr:sec:results}

\paragraph{Controlled primary result.}
\ControlledPrimaryResult

\paragraph{Safety and usefulness.}
\ControlledSafetyResult\ \ControlledRetentionResult

\ControlledMainResultTable
\ControlledMainResultFigure

\paragraph{Prior IID uncertainty-control result.}
In a disjoint 32-seed study, validation point selection violated 35/128
contracts (27.3\%), while the $\Gamma=1$ fixed-profile certificate violated
1/128 (0.8\%) and retained 52/102 oracle-safe opportunities (51.0\%). All four
dataset-level paired comparisons favored certification after Holm correction,
but the composite gate failed because GaitPDB's point violation rate, 5/32,
missed a registered 20\% baseline-severity floor. This supports uncertainty-
aware selection only; it is not evidence for $\Gamma>1$ robustness.

\paragraph{Negative results.}
\ControlledNegativeResults

\section{What Parts of VERA Matter?}
\label{iclr:sec:ablations}

The frozen ablation family isolates paired versus unpaired target loss;
balanced versus ordinary leakage; native versus retrained attackers; attacker
portfolio composition; IID versus shifted certification; scalar versus vector
eligibility; common versus anisotropic profiles; support checking;
multiplicity; evidence allocation; candidate screening and family size;
environment count; evidence size; exact versus generic bounds; and targeted
versus untargeted collection. Every ablation reuses registered objects or a
separate frozen stream and reports failed or null effects. None may replace a
failed primary endpoint.

\paragraph{Attacker stress.}
\ControlledHeldoutResult{} The held-out attacker provides stress evidence outside
the registered portfolio and is not part of the finite-sample certificate.
Failure against it limits portfolio scope; success does not imply resistance to
every learner.

\section{Evidence Growth and Failure Diagnosis}
\label{iclr:sec:diagnosis}

\ControlledAllocationResult{} The final analysis compares uniform allocation,
the locked square-score primary, and the additive program on disjoint common
random streams. It reports the continuous objective, deterministic integer
rounding gap, contract-level power surrogate, safety events, and safe retention.
The title is narrowed to ``Support-Aware'' because the registered usefulness
gate failed.

\ControlledEvidenceResult{} The inverse planning diagnostic reports continuous
and integer evidence requirements, per-cell top-ups, cap censoring, and the
provenance of every design-fold or oracle-only margin.

\ControlledGaitResult{} The GaitPDB analysis identifies limiting target or
attacker cells, current evidence, continuous and integer evidence requirements,
and recommended top-ups under fixed contracts. A top-up is a planning
diagnostic, not a post hoc threshold change and not a way to rescue a failed
confirmatory gate.

\section{Limitations, Responsible Use, and Conclusion}
\label{iclr:sec:limits}

VERA assumes bounded audit variables, independence between construction and
certification, a registered candidate and attacker family, and deployment
support covered by certification. Its density-ratio cap is declared rather than
learned from an unknown deployment. A finite attacker portfolio cannot certify
absence of all recoverable information. Target control applies to the
registered downstream learner and task, not every future use. Thresholds are
benchmark operating points. Algorithmic seeds measure fitting and split
variation but do not turn benchmark datasets into random deployment domains.
The controlled experiment is a finite-reference replay, and Camelyon17 is not
clinical validation.

The allocation result is restricted to separate target-environment and source-
class streams; a joint stratified design may reuse observations differently.
The support theorem is a boundary, not evidence that an unsupported cell is
unsafe. These distinctions are central to the claim rather than caveats added
after inspection.

\paragraph{Responsible use.}
Concept removal can hide a measured attribute without correcting structural
bias. A VERA certificate must not be represented as privacy, legal fairness, or
fitness for high-stakes deployment. It is a conditional decision statement:
this registered edit meets these paired target and attacker contracts for this
supported shift profile. When available evidence cannot justify that statement,
abstention is the scientifically correct output. The experiments use previously
released secondary benchmark data and recruit no new participants. The release
excludes raw biographies, comments, medical images, gait recordings,
participant-linked identifiers, and derivatives whose terms do not permit
redistribution.

\paragraph{AI assistance.}
OpenAI Codex was used extensively for ideation, literature discovery, theorem
and proof drafting, implementation, experiment orchestration, statistical code,
figures, manuscript drafting, and editing. No AI system is an author or a
citable source. Human authors are responsible for independently verifying every
claim, proof, implementation, reference, data right, and policy obligation
before submission.
